\section{OBJETIVOS DO CURSO}
\label{OBJETIVOS}

\subsection{Objetivo Geral}

O curso tem como objetivo formar engenheiros capazes de atuar na concepção, análise, desenvolvimento, integração e gestão de sistemas de controle e automação, fundamentando-se no desenvolvimento de competências profissionais que articulem conhecimentos científicos e tecnológicos, habilidades de aplicação prática e atitudes orientadas à responsabilidade técnica e social. Pretende-se preparar um profissional apto a atuar em contextos industriais dinâmicos, a assumir funções de coordenação e gestão de processos e equipes, bem como a desenvolver atividades de pesquisa aplicada e docência, caso opte pela continuidade de sua formação acadêmica.

\subsection{Objetivos Específicos}

A formação no \NOMECURSO organiza-se segundo o desenvolvimento integrado de competências, compreendidas em três dimensões interdependentes, a saber:

\begin{itemize}
	\item \textbf{Conhecimentos}: fundamentos de matemática, física, computação, eletrônica, controle, automação, gestão e produção industrial, que sustentam a compreensão dos princípios que regem a concepção e a operação de sistemas automatizados;
	\item \textbf{Habilidades}: capacidade de projetar, implementar, testar, manter e otimizar sistemas de controle e automação; analisar dados e indicadores de desempenho; aplicar normas técnicas; utilizar ferramentas digitais avançadas; comunicar-se com clareza em ambientes técnicos e acadêmicos; coordenar equipes e processos;
	\item \textbf{Atitudes}: postura ética, responsabilidade socioambiental, visão crítica e sistêmica, iniciativa, criatividade, liderança colaborativa, capacidade de investigação científica e comprometimento com a atualização profissional contínua.
\end{itemize}


Na Fig. \ref{fig:compet} é apresentada a interdependências entre os eixos que circundam a versão atual do PPC do \NOMECURSO.

\begin{figure}[h!]
\centering
		\caption{PPC baseado em competências}
		\begin{tikzpicture}[font=\sffamily, scale=0.7, every node/.style={transform shape}] 
			
			% Estilo dos blocos
			\tikzstyle{elemento} = [circle, draw=blue!60, fill=blue!10, thick, minimum size=3.2cm, text centered, align=center]
			\tikzstyle{ppc} = [rectangle, draw=orange!80, fill=orange!20, thick, rounded corners,
			minimum width=5cm, minimum height=1.2cm, text centered, align=center]
			
			% Coordenadas dos vértices do triângulo
			\node[elemento] (C) at (0,0) {Conhecimento};
			\node[elemento] (H) at (8,0) {Habilidade};
			\node[elemento] (A) at (4,6.9) {Atitudes};
			
			% PPC centralizado 
			\node[ppc, above=-6.5cm of A] (PPC) {\textbf{Foco em Competências}};
			
			% Conexões entre os elementos (arestas internas)
			%			\draw[very thick, blue!40] (C) -- (H) -- (A) -- (C);
			
			% Arcos externos conectando os círculos
			\draw[very thick, blue!70!black!50]
			let \p1 = (C), \p2 = (H), \p3 = (A) in
			% arco entre Conhecimento e Habilidade
			(C) .. controls ($(C)!0.5!(H) + (0,-2.2)$) .. (H)
			% arco entre Habilidade e Atitudes
			(H) .. controls ($(H)!0.5!(A) + (2,1.5)$) .. (A)
			% arco entre Atitudes e Conhecimento
			(A) .. controls ($(A)!0.5!(C) + (-2,1.5)$) .. (C);
			
			% Conexões do PPC com cada vértice
			\draw[{Latex[length=3mm]}-{Latex[length=3mm]}, orange!80, thick] (C) -- (PPC);
			
			\draw[{Latex[length=3mm]}-{Latex[length=3mm]}, orange!80, thick] (H) -- (PPC);
			\draw[{Latex[length=3mm]}-{Latex[length=3mm]}, orange!80, thick] (A) -- (PPC);
			
			% Textos explicativos
			\node[below=0.1cm of C, align=center, text width=4cm] {\footnotesize Base conceitual:\\ saber o quê e por quê?};
			\node[below=0.1cm of H, align=center, text width=4cm] {\footnotesize Base prática:\\ saber fazer};
			\node[above=0.1cm of A, align=center, text width=4cm] {\footnotesize Base ética e social:\\ saber ser};
			
		\end{tikzpicture}
		\caption*{Fonte: Autor, (2025)}
    \label{fig:compet}

\end{figure}


Com base nesses eixos, os objetivos específicos do  \NOMECURSO são:

\begin{itemize}
	\item Desenvolver competência científica e tecnológica para atuação qualificada no controle de processos e na automação industrial;
	\item Capacitar o estudante para atuar em equipes multidisciplinares e para exercer funções de gestão de projetos, processos e pessoas no contexto industrial;
	\item Formar profissionais aptos a avaliar e tomar decisões fundamentadas, considerando critérios técnicos, econômicos, ambientais e de segurança;
	\item Incentivar a investigação aplicada, o desenvolvimento de soluções inovadoras e a participação em projetos de pesquisa, extensão e inovação;
	\item Estimular a possibilidade de continuidade acadêmica em programas de pós-graduação, favorecendo o desenvolvimento de competências para atuação como pesquisador e docente;
	\item Favorecer a integração do estudante com o setor produtivo mediante estágios supervisionados, projetos integradores e cooperação com arranjos produtivos locais e regionais;
	\item Orientar para o cumprimento rigoroso das normas técnicas, regulamentações de segurança, saúde e meio ambiente;
	\item Promover postura empreendedora, liderança e capacidade de gestão orientada por responsabilidade social e sustentabilidade;
	\item Assegurar formação flexível, permitindo a incorporação de novas tecnologias emergentes no campo do controle e automação.
\end{itemize}

Este tripé está devidamente estruturado no PPC por meio de uma carga horária articulada que integra o Programa de Unidade Didática (PUD),  Projeto de Final de Curso (PFC), estágio curricular e curricularização, em consonância com os documentos-guia apresentados na Seção \ref{FUNDAMENTAÇÃO LEGAL}.